\section{Introduction}

The twelve faces of a regular dodecahedron occur in six opposite pairs.
Each pair determines an antipodal face-normal direction through the center
of the polyhedron. These six directions provide a natural coordinate system
for a restricted Euclidean realization family.

Fix the twelve face-normal directions and allow only the perpendicular
distances of the six opposite supporting plane pairs from the origin to
vary. Assigning one positive support number to each pair gives
\[
\mathbf h=(h_1,\ldots,h_6)\in\mathbb R_{>0}^6.
\]
The intersection of the corresponding symmetric slabs is an
origin-centered, centrally symmetric convex polytope.

The central question is:

\begin{quote}
How does the six-dimensional support register separate regular Euclidean
scale from relative shape variation?
\end{quote}

Define
\[
\bar h
=
\frac{1}{6}\sum_{i=1}^{6}h_i
\]
and
\[
\mathbf r
=
\mathbf h-\bar h\,\mathbf 1.
\]
Then
\[
\mathbf h
=
\bar h\,\mathbf 1+\mathbf r,
\qquad
\sum_{i=1}^{6}r_i=0.
\]

The algebraic decomposition is immediate. The main result of the paper is
its Euclidean interpretation.

Common multiplication of all support values scales the entire realization:
\[
P(c\mathbf h)=cP(\mathbf h).
\]
Conversely, while all defining planes remain active facets, a realization
homothetic to the regular reference dodecahedron must have six equal support
numbers. The uniform ray is therefore the unique regular scale direction,
and every nonzero residual gives a non-homothetic support variation.

The result is local to the connected chamber in support space that contains
the regular state and preserves dodecahedral incidence. Sufficiently small
support perturbations remain in this chamber by continuity and the strict
facet inequalities of the regular dodecahedron.

\subsection{Main results}

The paper establishes:

\begin{enumerate}
\item a complete six-coordinate description of the active fixed-normal,
origin-centered realization family;

\item the canonical orthogonal decomposition
\[
\mathbb R^6
=
\operatorname{span}\{\mathbf 1\}
\oplus
\left\{
\mathbf r\in\mathbb R^6:
\sum_i r_i=0
\right\};
\]

\item the identity
\[
P(c\mathbf h)=cP(\mathbf h),
\qquad
c>0;
\]

\item uniqueness of the regular scale direction while all defining planes
remain facet-defining;

\item local preservation of the dodecahedral combinatorial type; and

\item symmetry-compatible coordinates for mean scale, shape magnitude, and
normalized shape direction.
\end{enumerate}

\subsection{Scope}

All constructions take place in ordinary three-dimensional Euclidean space
and use standard convex polytope theory and finite-dimensional real linear
algebra.

The face-normal directions remain fixed, opposite supporting planes remain
equidistant from the origin, and central symmetry is preserved. Arbitrary
vertex displacement, independent face rotation, varying normal directions,
and non-centrally symmetric realizations are excluded.

No mechanics, elasticity, force, energy, dynamics, continuum model,
physical expansion, gravitation, or cosmological interpretation is
asserted.

\subsection{Organization}

Section~2 defines the support realization. Section~3 develops the canonical
split of the support register. Section~4 proves its Euclidean interpretation
as scale and shape. Section~5 records coordinate consequences, Section~6
gives examples, and Sections~7 and~8 discuss the scope and conclusions.
